\section{Prefix-Support Mismatch in Retained Harmful Suffixes}
\label{app:prefix-divergence}
This section formalizes the prefix-support mismatch mechanism identified in
\cref{sec:diagnosis}. We show why retaining \(\mathcal H_\kappa\)-tokens can be
unstable even when their local importance ratios are not tiny: the behavior
policy may assign non-negligible mass to these suffix tokens, while the current
policy assigns very small mass to the prefixes needed to reach them.

We follow the notation in Section~\ref{sec:preliminaries}: $\beta$ denotes the behavior policy, $\pi_\theta$ denotes the current policy, and $\rho_t(\theta)$ denotes the token-level importance ratio. For notational simplicity, we drop the response index and write a response as $y=(a_1,\ldots,a_T)$.  Assume a finite decoding horizon $T$; variable-length responses can be represented by padding with EOS tokens. Let $\lambda$ be any fixed distribution over token positions $t\in\{1,\ldots,T\}$ with $\lambda_t\ge 0$ and $\sum_t \lambda_t=1$.

For a token occurrence
\(z=(x,y,t)\), define the behavior token-occurrence measure
\[
    d\nu_\beta(x,y,t)
    =
    d\mathcal D(x)\,\lambda_t\,\beta(y\mid x),
\]
and the current-prefix occurrence measure
\[
    d\nu_\theta^{\mathrm{pre}}(x,y,t)
    =
    d\mathcal D(x)\,\lambda_t\,
    \pi_\theta(a_{<t}\mid x)\,
    \beta(a_{t:T}\mid x,a_{<t}).
\]
Thus \(\nu_\theta^{\mathrm{pre}}\) changes only the prefix occupancy from
\(\beta\) to \(\pi_\theta\), while keeping the replay suffix conditional under
\(\beta\). Whenever \(\beta(a_{<t}\mid x)>0\),
\[
    \frac{d\nu_\theta^{\mathrm{pre}}}{d\nu_\beta}(x,y,t)
    =
    R_{t-1}(\theta)
    :=
    \frac{\pi_\theta(a_{<t}\mid x)}{\beta(a_{<t}\mid x)}
    =
    \prod_{j=1}^{t-1}
    \rho_j(\theta),
\]
where
\[
    \rho_t(\theta)
    :=
    \frac{\pi_\theta(a_t\mid x,a_{<t})}
         {\beta(a_t\mid x,a_{<t})}.
\]

Recall from Section~\ref{sec:diagnosis} that the {trigger boundary} is the first trigger-token position
\[
    \kappa(y)
    :=
    \min\{t:\, A(y)<0,\ \rho_t(\theta)<\tau_c\},
\]
with $\kappa(y)=\infty$ if the set is empty. The non-trigger suffix is
\[
    \mathcal H_\kappa(y)
    :=
    \{t>\kappa(y):\, A(y)<0,\ \rho_t(\theta)\ge \tau_c\}.
\]
Let \(\mathcal U_{\mathcal H}(\theta)\subseteq\{(x,y,t): t\in \mathcal H_\kappa(y)\}\) denote any
retained subset of these \(\mathcal H_\kappa\)-tokens. If the update rule retains all of \(\mathcal H_\kappa(y)\),
then \(\mathcal U_{\mathcal H}(\theta)=\{(x,y,t):t\in \mathcal H_\kappa(y)\}\).

We use the Pearson chi-square convention
\[
    \chi^2(P\|Q)
    :=
    \int \left(\frac{dP}{dQ}-1\right)^2\,dQ,
\]
with \(\chi^2(P\|Q)=+\infty\) when \(P\) is not absolutely continuous with
respect to \(Q\).

\begin{theorem}[Retained \(\mathcal H_\kappa\)-tokens can induce arbitrarily large prefix chi-square divergence]
\label{thm:retained-H-chi2}
Let
\[
    m_{\mathcal H}
    :=
    \nu_\beta(\mathcal U_{\mathcal H}),
    \qquad
    q_{\mathcal H}
    :=
    \nu_\theta^{\mathrm{pre}}(\mathcal U_{\mathcal H})
    =
    \mathbb E_{\nu_\beta}
    \left[
        R_{t-1}(\theta)\mathbf 1\{(x,y,t)\in\mathcal U_{\mathcal H}\}
    \right].
\]
Assume \(m_{\mathcal H}>0\) and
\[
    q_{\mathcal H} \le r\,m_{\mathcal H}
    \qquad\text{for some } r\in(0,1).
\]
Then
\[
    \chi^2\!\left(\nu_\beta\middle\|\nu_\theta^{\mathrm{pre}}\right)
    \ge
    \chi^2\!\left(
        \mathrm{Bern}(m_{\mathcal H})
        \middle\|
        \mathrm{Bern}(q_{\mathcal H})
    \right)
    =
    \frac{(m_{\mathcal H}-q_{\mathcal H})^2}{q_{\mathcal H}(1-q_{\mathcal H})}.
\]
Consequently,
\[
    \chi^2\!\left(\nu_\beta\middle\|\nu_\theta^{\mathrm{pre}}\right)
    \ge
    \frac{m_{\mathcal H}(1-r)^2}{r}.
\]
In particular, for fixed \(m_{\mathcal H}>0\), the lower bound diverges as
\(r\downarrow 0\). If \(q_{\mathcal H}=0<m_{\mathcal H}\), then
\[
    \chi^2\!\left(\nu_\beta\middle\|\nu_\theta^{\mathrm{pre}}\right)
    =
    +\infty .
\]
\end{theorem}

\begin{proof}
Let \(P=\nu_\beta\), \(Q=\nu_\theta^{\mathrm{pre}}\), and
\(U=\mathcal U_{\mathcal H}\). If \(Q(U)=0<P(U)\), then \(P\) is not absolutely continuous
with respect to \(Q\), so \(\chi^2(P\|Q)=+\infty\).

Now assume \(q_{\mathcal H}=Q(U)>0\). Let \(L=dP/dQ\). By Cauchy--Schwarz,
\[
    \int_U L^2\,dQ
    \ge
    \frac{\left(\int_U L\,dQ\right)^2}{Q(U)}
    =
    \frac{P(U)^2}{Q(U)}
    =
    \frac{m_{\mathcal H}^2}{q_{\mathcal H}}.
\]
Applying the same argument to \(U^c\),
\[
    \int_{U^c} L^2\,dQ
    \ge
    \frac{P(U^c)^2}{Q(U^c)}
    =
    \frac{(1-m_{\mathcal H})^2}{1-q_{\mathcal H}}.
\]
Therefore,
\[
\begin{aligned}
    \chi^2(P\|Q)
    &=
    \int L^2\,dQ -1 \\
    &\ge
    \frac{m_{\mathcal H}^2}{q_{\mathcal H}}
    +
    \frac{(1-m_{\mathcal H})^2}{1-q_{\mathcal H}}
    -1 \\
    &=
    \frac{(m_{\mathcal H}-q_{\mathcal H})^2}{q_{\mathcal H}(1-q_{\mathcal H})}.
\end{aligned}
\]
Using \(q_{\mathcal H}\le r m_{\mathcal H}\), we have
\[
    m_{\mathcal H}-q_{\mathcal H} \ge m_{\mathcal H}(1-r),
    \qquad
    q_{\mathcal H}(1-q_{\mathcal H})\le q_{\mathcal H} \le r m_{\mathcal H}.
\]
Thus
\[
    \chi^2(P\|Q)
    \ge
    \frac{m_{\mathcal H}^2(1-r)^2}{r m_{\mathcal H}}
    =
    \frac{m_{\mathcal H}(1-r)^2}{r}.
\]
This proves the claim.
\end{proof}

\begin{corollary}[The blow-up is compatible with non-tiny local ratios]
\label{cor:non-tiny-local-ratio}
Fix any threshold \(\tau_c\in(0,1]\) and any target \(M>0\). There exist
behavior and current policies \(\beta\) and \(\pi_\theta\), and a retained
harmful suffix set \(\mathcal U_{\mathcal H}\), such that every retained suffix token in
\(\mathcal U_{\mathcal H}\) satisfies
\[
    \rho_t(\theta)\ge \tau_c,
\]
but
\[
    \chi^2\!\left(\nu_\beta\middle\|\nu_\theta^{\mathrm{pre}}\right)
    >
    M.
\]
\end{corollary}

\begin{proof}
Consider one prompt and two token positions. Let the first-token vocabulary
contain two tokens \(c\) and \(b\). Choose \(m\in(0,1)\), and choose
\[
    r < \min\{\tau_c,\,1/2,\,\lambda_2 m/(4M)\},
\]
where \(\lambda_2>0\) is the position weight assigned to the second token.

Let
\[
    \beta(c\mid x)=m,
    \qquad
    \pi_\theta(c\mid x)=rm,
\]
with the remaining probability assigned to \(b\). After prefix \(c\), let both
policies deterministically emit the second token \(h\):
\[
    \beta(h\mid x,c)=1,
    \qquad
    \pi_\theta(h\mid x,c)=1.
\]
Assign negative advantage to the trajectory \(y=(c,h)\), i.e.,
\(A(y)<0\).

For \(y=(c,h)\), the first token has
\[
    \rho_1(\theta)
    =
    \frac{\pi_\theta(c\mid x)}{\beta(c\mid x)}
    =
    r
    <
    \tau_c,
\]
so \(\kappa(y)=1\). The second token has
\[
    \rho_2(\theta)
    =
    \frac{\pi_\theta(h\mid x,c)}{\beta(h\mid x,c)}
    =
    1
    \ge
    \tau_c,
\]
so \(2\in \mathcal H_\kappa(y)\). Its prefix occupancy is
\[
    R_1(\theta)=\rho_1(\theta)=r.
\]
Let the update rule retain this second-token occurrence, and define
\[
    \mathcal U_{\mathcal H}=\{(x,(c,h),2)\}.
\]
Then
\[
    m_{\mathcal H}=\nu_\beta(\mathcal U_{\mathcal H})=\lambda_2 m,
    \qquad
    q_{\mathcal H}=\nu_\theta^{\mathrm{pre}}(\mathcal U_{\mathcal H})=\lambda_2 r m.
\]
By Theorem~\ref{thm:retained-H-chi2},
\[
    \chi^2\!\left(\nu_\beta\middle\|\nu_\theta^{\mathrm{pre}}\right)
    \ge
    \frac{\lambda_2 m(1-r)^2}{r}.
\]
Since \(r<1/2\), \((1-r)^2>1/4\), and therefore
\[
    \chi^2\!\left(\nu_\beta\middle\|\nu_\theta^{\mathrm{pre}}\right)
    >
    \frac{\lambda_2 m}{4r}
    >
    M.
\]
Thus the retained token can have a perfectly benign local ratio
\(\rho_2(\theta)=1\), while the induced prefix mismatch is arbitrarily large.
\end{proof}

\begin{remark}[Direction of the Pearson divergence]
\label{rem:chi2-direction}
The divergence in Theorem~\ref{thm:retained-H-chi2} is the reverse Pearson
divergence \(\chi^2(\nu_\beta\|\nu_\theta^{\mathrm{pre}})\), which is sensitive
to behavior-supported prefixes that have very small current occupancy. The
opposite direction satisfies
\[
    \chi^2\!\left(\nu_\theta^{\mathrm{pre}}\middle\|\nu_\beta\right)
    =
    \mathbb E_{\nu_\beta}
    \left[
        \left(R_{t-1}(\theta)-1\right)^2
    \right],
\]
so \(R_{t-1}\ll 1\) on a set of fixed behavior mass does not by itself force an
unbounded contribution in that direction. This is why the missing-support
statement should be phrased as a behavior-to-current, or reverse-Pearson,
prefix divergence.
\end{remark}